% Generated draft: requires mathematical and source review.
\documentclass[11pt]{article}
\usepackage[margin=25mm]{geometry}
\usepackage{fontspec}
\setmainfont{Latin Modern Roman}
\usepackage{amsmath,amssymb,mathtools,unicode-math}
\setmathfont{Latin Modern Math}
\usepackage{enumitem,needspace,xurl,hyperref}
\setlength{\parindent}{0pt}
\setlength{\parskip}{5pt}
\setlength{\emergencystretch}{3em}
\newcommand{\N}{\mathbb{N}}
\newcommand{\Z}{\mathbb{Z}}
\newcommand{\Q}{\mathbb{Q}}
\newcommand{\R}{\mathbb{R}}
\newcommand{\C}{\mathbb{C}}
\newcommand{\F}{\mathbb{F}}
\newcommand{\A}{\mathbb{A}}
\newcommand{\PP}{\mathbb{P}}
\newcommand{\E}{\mathbb{E}}
\newcommand{\eps}{\varepsilon}
\newcommand{\dd}{\,\mathrm{d}}
\begin{document}
\section*{UnsolvedMath ID: 20002224}
\textbf{Source problem code:} catalogue entry
\subsection{Definitions and mathematical statement}
Let $G$ be a finite simple graph with $n\ge1$ vertices, $m$ edges and degrees $d(v)$. Its combinatorial Laplacian is $L=D-A$, with eigenvalues $\lambda_1\ge\cdots\ge\lambda_n\ge0$. Define the conjugate degree sequence, padded to length $n$, by
\[
 d_i^*=\#\{v:d(v)\ge i\},\qquad1\le i\le n.
\]
Both sequences sum to $2m$. Majorization $\lambda\prec d^*$ means
\[
 \sum_{i=1}^k\lambda_i\le\sum_{i=1}^kd_i^*
      =\sum_v\min(d(v),k)\quad(1\le k<n),
 \qquad\sum_{i=1}^n\lambda_i=\sum_{i=1}^nd_i^*.
\]
Put $\mu=2m/n$ and $E_L(G)=\sum_i|\lambda_i-\mu|$.

\textbf{Question 6, Grone--Merris conjecture.} Is $\lambda\prec d^*$? The source reports the workshop consequence
\[
 E_L(G)\le B(G):=\sum_{i=1}^n|d_i^*-\mu|.
\]
The supplied aid attributes an affirmative theorem for every finite simple graph to Hua Bai (2011). This is a named literature result as supplied, not a new proof by this draft.

\textbf{Supplied exact energy and equality refinement.} The degree-only bound can be written
\[
 B(G)=2\max_{0\le k\le n}
       \left(\sum_v\min(d(v),k)-\frac{2mk}{n}\right).
\]
For a decreasing sequence $x$ of mean $\mu$, define its mean-crossing index set
\[
 I(x)=\{k\in\{0,\ldots,n\}:\#\{i:x_i>\mu\}\le k
                              \le\#\{i:x_i\ge\mu\}\}.
\]
The aid asserts that equality $E_L(G)=B(G)$ holds precisely when some $k\in I(\lambda)\cap I(d^*)$ has a tight Grone--Merris prefix,
\[
 \sum_{i=1}^k\lambda_i=\sum_{i=1}^kd_i^*.
\]
The inclusive mean-crossing convention handles repeated entries equal to $\mu$, and $k=0,n$ handles the edgeless/degenerate case without introducing undefined indices.

For the complete bipartite graph $K_{a,b}$, $1\le a\le b$, the aid gives the exact deficit
\[
 B(K_{a,b})-E_L(K_{a,b})=2a(a-1).
\]
Hence equality in this family occurs exactly for $a=1$, the stars. This is a family-specific equality conclusion, not a claim that only stars attain equality among all simple graphs.
\subsection{Short English statement}
The source asks whether Laplacian eigenvalues are majorized by the conjugate degree sequence; the aid cites Bai's affirmative theorem. Its refinements give a degree-only Laplacian-energy bound, an exact tight-prefix equality test, and a complete-bipartite gap showing equality there exactly for stars.
\subsection{Sources}
\begin{itemize}
\item Local UnsolvedMath registry, permanent ID 20002224, source code catalogue entry; canonical \texttt{aim-linear-algebra-notes.json}, notes[17], statement status reconstructed unverified.
\item AIM \emph{Spectra of families of matrices described by graphs, digraphs, and sign patterns}, Question 6 and workshop energy observation: \url{https://aimath.org/WWN/matrixspectrum/matrixspectrum.pdf}.
\item Grone--Merris conjecture; the supplied aid attributes its proof to Hua Bai (2011), without a full paper title, publication details or theorem label.
\item Registry research attempt 1, dated 2026-08-14: full solution, problem status at run solved, rigor ``proved'', novelty confidence low. The explicitly machine-generated, non-peer-reviewed aid requires independent verification before citation.
\end{itemize}
\end{document}
